\section{Arithmetic targets and the remaining mathematical work}
\label{sec:arithmetic-targets}

\subsection{A concrete nonabelian test case}

Let $F=\mathbb Q(\alpha)$ with $\alpha^3-\alpha-1=0$.
The polynomial has no rational root and discriminant $-23$, so its
splitting field has Galois group $S_3$. The standard two-dimensional
character $\chi$ has values $2,0,-1$ on the identity, transpositions,
and three-cycles. The permutation character on the three roots is
$1+\chi$. Thus its unramified local data are
\[
\begin{array}{c|ccc}
\text{factor degrees modulo }p & 1+1+1&1+2&3\\ \hline
\text{fixed roots} &3&1&0\\
\chi(\operatorname{Frob}_p)&2&0&-1
\end{array}
\]
This supplies a definite arithmetic target: control the nonabelian
factor in the Euler-factor identity
$\zeta_F(s)=\zeta(s)L(s,\chi)$.

For precision, define the nonnegative coefficients $a_F(n)$ by
\[
 -\frac{\zeta'_F}{\zeta_F}(s)=\sum_{n\ge1}\frac{a_F(n)}{n^s}
       \qquad(\operatorname{Re}s>1).
\]
Equivalently, sum $\log N\mathfrak p$ over prime-ideal powers
$\mathfrak p^r$ of norm $n$. For $p\ne23$,
$a_F(p^r)=\#\operatorname{Fix}(\operatorname{Frob}_p^r)\log p$.
Define
\[
 R_F(x)=\sum_{n\le x}a_F(n)\left(\frac1n-\frac1x\right)
                    -\log x-c_F,
\]
where $c_F=-1-\gamma_F$ and $\gamma_F$ is the constant term of
$\zeta'_F/\zeta_F(s)+1/(s-1)$ at $s=1$.
This mirrors the source's weighted remainder in
\src{New/Reduction/Zeros/PrimeError.lean}{PrimeError}, while retaining
the new Frobenius information.

A clean sufficient goal is the two-sided estimate
\begin{equation}
 \label{eq:cubic-remainder-goal}
 |R_F(x)|\ll_\varepsilon x^{-1/2+\varepsilon}
                  \qquad\text{for every }\varepsilon>0.
\end{equation}
Indeed, termwise integration in the initial half-plane gives
\[
 \int_1^\infty R_F(x)x^{-s}\,dx
 =\frac{-\zeta'_F/\zeta_F(s)}{s(s-1)}
       -\frac1{(s-1)^2}-\frac{c_F}{s-1}.
\]
The estimate would make the integral holomorphic for
$\operatorname{Re}s>1/2$, so the logarithmic derivative could have no
pole there except the prescribed one at $1$. This excludes zeros in
that half-plane; the functional equation supplies the usual symmetry.
This is a proved analytic reduction, \emph{not} a proof of
\eqref{eq:cubic-remainder-goal}.

The source achieves a one-sided bound in its particular setting
(\src{New/Reduction/OneSidedBound.lean}{OneSidedBound}).
Replacing its congruence weights by the table above does not transfer
the proof. A new construction must realize these weights as local
divisibility gains, retain almost-full spanning with controlled index,
and produce a determinant lower bound with the same main term.
Those are the unresolved arithmetic steps. Reusing a weaker one-sided
criterion would additionally require its pole and cancellation
hypotheses to be checked for the chosen factors.

\subsection{Cyclotomic and parity restrictions}

The explicit odd-cocycle bound
\texttt{exists\_length\_shaOdd\_le} in
\src{Lib/Cyclotomic/OddSha.lean}{OddSha}
assumes a CM field $E\subseteq\mathbb Q(\mu_N)$. For a fixed twist range
$0\le e<t$, it gives a nonzero integer
\[
 X\le 2^{4(N+t+3)^{27}}h_E
\]
whose valuations bound the specified groups at every torsion precision.
It does not supply a conductor-independent bound as $E$ varies, nor an
immediate replacement over a nonabelian base.

The word ``odd'' also requires translation. For $e\ne0$, the source
imposes $z(cgc)=-z(g)$ with coefficients in a Tate twist $(-e)$.
The coefficient action of complex conjugation contributes $(-1)^e$,
so the usual cohomological sign is $(-1)^{e+1}$.
The case $e=0$ has no such parity condition; see
\src{Lib/GaloisRep/Main/SelmerModules.lean}{the definition of shaOdd}.
The denominator discussion explicitly says the even classes relevant
to Vandiver do not enter
(\src{New/Congruences/FirstOrder/DenominatorBound.lean}{DenominatorBound}).
Consequently neither Vandiver nor a new algebraic $K$-group vanishing
follows by matching names or parities informally.

\subsection{Priorities after the completed extensions}

The following ordering reflects the amount of missing structure.

\begin{enumerate}
\item \textbf{Formalize the isolated algebraic interfaces.}
Extract the determinant comparison and the finite-level stopping lemma
with their complete hypotheses. Their proofs and source interfaces are
already available, making this the most tractable verification task.
\item \textbf{Verify one concrete arithmetic Selmer datum.}
Use the general twisting theorem, rather than its specialized
automorphic endpoint. Specify the representation, scalar inertia,
local conditions, and all terms of $\beta$. The source's
\texttt{congruenceTerm\_le} gives a useful certificate: a common scalar
witness with unit difference forces that global defect to vanish
(\src{New/Congruences/Selmer/SelmerData.lean}{SelmerData}).
An exact comparison to the intended arithmetic group remains mandatory.
\item \textbf{Construct the cubic Frobenius divisibility detector.}
Start with the three-entry table above and an explicit integral object.
Prove its local valuation formula before attempting the global
spanning and norm estimates. These quantitative estimates are the
substantial obstacle to a new nonabelian zero theorem.
\end{enumerate}

This report does not resolve a new major arithmetic conjecture.
The completed graph and specialization results are separate mathematical
advances; the arithmetic analysis specifies where additional ideas are
needed and prevents the available formal theorems from being applied
outside their hypotheses.
